%
% Copyright (C) 1995, 2000, 2013 American Mathematical Society.
% Copyright (C) 2016-2025 LaTeX Project and American Mathematical Society.
%
% This work may be distributed and/or modified under the
% conditions of the LaTeX Project Public License, either version 1.3c
% of this license or (at your option) any later version.
% The latest version of this license is in
%   https://www.latex-project.org/lppl.txt
% and version 1.3c or later is part of all distributions of LaTeX
% version 2005/12/01 or later.
%
% This work has the LPPL maintenance status `maintained'.
%
% The Current Maintainer of this work is the LaTeX Project.
%
\documentclass[leqno,titlepage,openany]{amsldoc}[1999/12/13]

\def\MaintainedByLaTeXTeam#1{%
\begin{center}%
\fbox{\fbox{\begin{tabular}{@{}l@{}}%
This file is maintained by the \LaTeX{} Project team.\\%
Bug reports can be opened (category \texttt{#1}) at\\%
\url{https://latex-project.org/bugs/}.\end{tabular}}}\end{center}}

\ifx\UndEfiNed\url
  \ClassError{amsldoc}{%
    This version of amsldoc.tex must be processed\MessageBreak
    with a newer version of amsldoc.cls (2.02 or later).}{}
\fi

\renewcommand\bibname{References}

%    Definitions that should be moved to amsldoc.cls
%    when it is next updated
\makeatletter
\let\cleardouble@page\cleardoublepage
\AtBeginDocument{%
  \ifx\cleardouble@page\cleardoublepage
    \def\cleardoublepage{\clearpage{\pagestyle{empty}\cleardouble@page}}
  \fi
}

% #1 is empty or one of class, environment, option, package
% could use different font in each case
\makeatletter
\def\@category@index#1#2{%
    \autoindex{%
        #2\actualchar\string\index#1font{#2}%
        \ifx\@nil#1\@nil\else\space#1\fi
    }%
}

\let\indexpackagefont\textsf
\let\indexclassfont\textsf
\let\indexoptionfont\texttt
\let\indexenvironmentfont\texttt
\let\indexfont\texttt
\makeatother
%    For a package that won't be indexed, i.e. amsmath
\newcommand{\nipkg}{\textsf}
\def\tikz/{Ti\textit{k}Z}
%    End definitions for amsldoc.cls

\title{User's Guide for the \nipkg{amsmath} Package (Version~2.1)}
\author{American Mathematical Society, \LaTeX\ Project}
\date{1999-12-13\\(revised 2002-02-25, 2016-11-14, 2018-04-05, 2019-10-14, 2020-02-18)}
\makeatletter
\def\@thanks{\bigskip\MaintainedByLaTeXTeam{amslatex}}
\makeatother

% Very minor changes are not listed in the \date above:
%
% 2021-08-22: reported typo fixed mathools -> mathtools
% 2025-05-23: drop spurious option list


%    Use the amsmath package and amscd package in order to print
%    examples.
\usepackage{amsmath}
\usepackage{amscd}

\usepackage{imakeidx}
%    Need alternate coding for vertical bar to get it into index.
%    Based on \qcbang and \cnbang in amsldoc.cls, refined by David Carlisle.
%    Give it a sort field that will be at the end of the symbols;
%    must avoid " and | which have special meanings.
\DeclareRobustCommand{\qcvert}{\qc\|\index{{z1}@{\ntt\char'174}}}
\DeclareRobustCommand{\cnvert}{%
  \ncn{\|}\index{{z0}@{\ntt\bslchar\char'174}}}
%    Need alternate index sort field for bang to allow & to sort first
\DeclareRobustCommand{\altcnbang}{%
  \ncn{\!}\index{"!b@{\ntt\bslchar\qcbang}}}
\DeclareRobustCommand{\cnamp}{%
  \index{"!a@\texttt{\&}}}

\makeindex[intoc] % generate index data
\providecommand{\see}[2]{\textit{see} #1}

%    The amsldoc class includes a number of features useful for
%    documentation about TeX, including:
%
%    ---Commands \tex/, \amstex/, \latex/, ... for uniform treatment
%    of the various logos and easy handling of following spaces.
%
%    ---Commands for printing various common elements: \cn for command
%    names, \fn for file names (including font-file names), \env for
%    environments, \pkg and \cls for packages and classes, etc.

%    Many of the command names used here are rather long and will
%    contribute to poor linebreaking if we follow the \latex/ practice
%    of not hyphenating anything set in tt font; instead we selectively
%    allow some hyphenation.
\allowtthyphens % defined in amsldoc.cls

\hyphenation{ac-cent-ed-sym-bol add-to-counter add-to-length align-at
  aligned-at allow-dis-play-breaks ams-art ams-cd ams-la-tex amsl-doc
  ams-symb ams-tex ams-text ams-xtra bmatrix bold-sym-bol cen-ter-tags
  eqn-ar-ray idots-int int-lim-its latex med-space neg-med-space
  neg-thick-space neg-thin-space no-int-lim-its no-name-lim-its
  over-left-arrow over-left-right-arrow over-right-arrow pmatrix
  qed-sym-bol set-length side-set small-er tbinom the-equa-tion
  thick-space thin-space un-der-left-arrow un-der-left-right-arrow
  un-der-right-arrow use-pack-age var-inj-lim var-proj-lim vmatrix
  xalign-at xx-align-at}

%    Prepare for illustrating the \vec example
\newcommand{\vect}[1]{\mathbf{#1}}

\newcommand{\booktitle}[1]{\textit{#1}}
\newcommand{\journalname}[1]{\textit{#1}}
\newcommand{\seriesname}[1]{\textit{#1}}

%    Command to insert and index a particular phrase. Doesn't work for
%    certain kinds of special characters in the argument.
\newcommand{\ii}[1]{#1\index{#1}}

\newcommand{\vstrut}[1]{\vrule width0pt height#1\relax}

%    An environment for presenting comprehensive address information:
\newenvironment{infoaddress}{%
  \par\topsep\medskipamount
  \trivlist\centering
  \item[]%
  \begin{minipage}{.7\columnwidth}%
  \raggedright
}{%
  \end{minipage}%
  \endtrivlist
}

\newenvironment{eqxample}{%
  \par\addvspace\medskipamount
  \noindent\begin{minipage}{.5\columnwidth}%
  \def\producing{\end{minipage}\begin{minipage}{.5\columnwidth}%
    \hbox\bgroup\kern-.2pt\vrule width.2pt%
      \vbox\bgroup\parindent0pt\relax
%    The 3pt is to cancel the -\lineskip from \displ@y
    \abovedisplayskip3pt \abovedisplayshortskip\abovedisplayskip
    \belowdisplayskip0pt \belowdisplayshortskip\belowdisplayskip
    \noindent}
}{%
  \par
%    Ensure that a lonely \[\] structure doesn't take up width less than
%    \hsize.
  \hrule height0pt width\hsize
  \egroup\vrule width.2pt\kern-.2pt\egroup
  \end{minipage}%
  \par\addvspace\medskipamount
}

%    Errors in output shouldn't be \texttt, and they may also be longer
%    than one line, but they should be identified by the "errorbullet".
%    (They were formerly labeled as aubsections.)
\newenvironment{erroro}[1]{%
  \par\addvspace\medskipamount
  \noindent\hangindent\parindent
  \hbox to\parindent{\errorbullet\hfil}\ignorespaces
  #1\par\vspace{\smallskipamount}
}{%
  \par\addvspace\medskipamount
}

%    The chapters are so short, perhaps we shouldn't call them by the
%    name `Chapter'. We make \chaptername read an argument in order to
%    remove a following \space or "{} " (both possibilities are present
%    in book.cls).

\renewcommand{\chaptername}[1]{}
\newcommand{\chapnum}[1]{\mdash #1\mdash }
\makeatletter
\def\@makechapterhead#1{%
  \vspace{1.5\baselineskip}%
  {\parindent \z@ \raggedright \reset@font
    \ifnum \c@secnumdepth >\m@ne
      \large\bfseries \chapnum{\thechapter}%
      \par\nobreak
      \vskip.5\baselineskip\relax
    \fi
    #1\par\nobreak
    \vskip\baselineskip
  }}
\makeatother

%    A command for ragged-right parbox in a tabular.
\newcommand{\rp}{\let\PBS\\\raggedright\let\\\PBS}

%    Non-indexed file name
\newcommand{\nfn}[1]{\texttt{#1}}

%    For the examples in the math spacing table.
%%\newcommand{\lspx}{\mbox{\rule{5pt}{.6pt}\rule{.6pt}{6pt}}}
%%\newcommand{\rspx}{\mbox{\rule[-1pt]{.6pt}{7pt}%
%%  \rule[-1pt]{5pt}{.6pt}}}
\newcommand{\lspx}{\mathord{\Rightarrow\mkern-1mu}}
\newcommand{\rspx}{\mathord{\mkern-1mu\Leftarrow}}
\newcommand{\spx}[1]{$\lspx #1\rspx$}

%    For a list of characters representing document input.
\newcommand{\clist}[1]{%
  \mbox{\ntt\spaceskip.2em plus.1em \xspaceskip\spaceskip#1}}

%    Fix weird \latex/ definition of rightmark.
\makeatletter
\def\rightmark{\expandafter\@rightmark\botmark{}{}}
%    Also turn off section marks.
\let\sectionmark\@gobble
\renewcommand{\chaptermark}[1]{%
  \uppercase{\markboth{\rhcn#1}{\rhcn#1}}}
\newcommand{\rhcn}{\thechapter. }
\makeatother

%    Include down to \section but not \subsection, in toc:
\setcounter{tocdepth}{1}

%    Subheadings for the bibliography
\newif\iffirstbibsubhead \firstbibsubheadtrue
\newcommand{\bibsubhead}[1]{%
  \iffirstbibsubhead \firstbibsubheadfalse
  \else \addvspace{\medskipamount}
  \fi
  \item[]\hspace*{-\leftmargin}\textbf{#1}\par
  \vspace{\smallskipamount}
}

\DeclareMathOperator{\ix}{ix}
\DeclareMathOperator{\nul}{nul}
\DeclareMathOperator{\End}{End}
\DeclareMathOperator{\xxx}{xxx}

\usepackage{hyperref}
\hypersetup{pdfborder=0 0 0}
\makeatletter
\let\oldcs\cs
\def\cs#1{\texorpdfstring{\oldcs{#1}}{\@backslashchar\@backslashchar#1}}
\let\cn\cs
\makeatother
\begin{document}
\section{Continued fractions}

The continued fraction\index{continued fractions}
\begin{equation}
\cfrac{1}{\sqrt{2}+
 \cfrac{1}{\sqrt{2}+
  \cfrac{1}{\sqrt{2}+\cdots
}}}
\end{equation}
can be obtained by typing
{\samepage
\begin{verbatim}
\cfrac{1}{\sqrt{2}+
 \cfrac{1}{\sqrt{2}+
  \cfrac{1}{\sqrt{2}+\dotsb
}}}
\end{verbatim}
}% End of \samepage
This produces better-looking results than straightforward use of
\cn{frac}. Left or right placement of any of the numerators is
accomplished by using \cn{cfrac}|[l]| or \cn{cfrac}|[r]| instead of
\cn{cfrac}.\index{options!adjust positioning}%
\index{fractions|)}

%% ------------------------------------------------------------------ %%

\section{Smash options}

The command \cn{smash} is used to typeset a subformula with an
effective height and depth of zero, which is sometimes
useful in adjusting the subformula's position with respect to adjacent
symbols. With the \nipkg{amsmath} package \cn{smash} has optional
arguments\index{options!behavior of particular options}
|[t]|\index{t option@\texttt{t} (top) option} and
|[b]|\index{b option@\texttt{b} (bottom) option} because occasionally it is
advantageous to be able to \qq{smash} only the top or only the bottom
of something while retaining the natural depth or height. For example,
when adjacent radical symbols are unevenly sized or positioned because
of differences in the height and depth of their contents, \cn{smash}
can be employed to make them more consistent. Compare
$\sqrt{x}+\sqrt{y}+\sqrt{z}$ and $\sqrt{x}+\sqrt{\smash[b]{y}}+\sqrt{z}$,
where the latter was produced by
\verb"$\sqrt{x}" \verb"+"
\verb"\sqrt{"\5\verb"\smash[b]{y}}" \verb"+" \verb"\sqrt{z}$".

%% ------------------------------------------------------------------ %%

\section{Delimiters}

\index{delimiters|(}

\subsection{Delimiter sizes}\label{bigdel}

The automatic delimiter sizing done by \cn{left} and \cn{right} has two
limitations: First, it is applied mechanically to produce delimiters
large enough to encompass the largest contained item, and second, the
range of sizes is not even approximately continuous but has fairly large
quantum jumps. This means that a math fragment that is infinitesimally
too large for a given delimiter size will get the next larger size, a
jump of 3pt or so in normal-sized text. There are two or three
situations where the delimiter size is commonly adjusted, using a set of
commands that have `big' in their names.%
\index{delimiters!fixed size|(}
\begin{ctab}{l|llllll}
Delimiter&
  text& \ncn{left}& \ncn{bigl}& \ncn{Bigl}& \ncn{biggl}& \ncn{Biggl}\\
size&
  size& \ncn{right}& \ncn{bigr}& \ncn{Bigr}& \ncn{biggr}& \ncn{Biggr}\\
\hline
Result\vstrut{5ex}&
  $\displaystyle(b)(\frac{c}{d})$&
  $\displaystyle\left(b\right)\left(\frac{c}{d}\right)$&
  $\displaystyle\bigl(b\bigr)\bigl(\frac{c}{d}\bigr)$&
  $\displaystyle\Bigl(b\Bigr)\Bigl(\frac{c}{d}\Bigr)$&
  $\displaystyle\biggl(b\biggr)\biggl(\frac{c}{d}\biggr)$&
  $\displaystyle\Biggl(b\Biggr)\Biggl(\frac{c}{d}\Biggr)$
\end{ctab}
The first kind of situation is a cumulative operator with limits above
and below. With \cn{left} and \cn{right} the delimiters usually turn out
larger than necessary, and using the |Big| or |bigg|
sizes\index{big@\cn{big}, \cn{Big}, \cn{bigg}, \dots\ delimiters}
instead gives better results:
\begin{equation*}
\left[\sum_i a_i\left\lvert\sum_j x_{ij}\right\rvert^p\right]^{1/p}
\quad\text{versus}\quad
\biggl[\sum_i a_i\Bigl\lvert\sum_j x_{ij}\Bigr\rvert^p\biggr]^{1/p}
\end{equation*}
\begin{verbatim}
\biggl[\sum_i a_i\Bigl\lvert\sum_j x_{ij}\Bigr\rvert^p\biggr]^{1/p}
\end{verbatim}
The second kind of situation is clustered pairs of delimiters where
\cn{left} and \cn{right} make them all the same size (because that is
adequate to cover the encompassed material) but what you really want
is to make some of the delimiters slightly larger to make the nesting
easier to see.
\begin{equation*}
\left((a_1 b_1) - (a_2 b_2)\right)
\left((a_2 b_1) + (a_1 b_2)\right)
\quad\text{versus}\quad
\bigl((a_1 b_1) - (a_2 b_2)\bigr)
\bigl((a_2 b_1) + (a_1 b_2)\bigr)
\end{equation*}
\begin{verbatim}
\left((a_1 b_1) - (a_2 b_2)\right)
\left((a_2 b_1) + (a_1 b_2)\right)
\quad\text{versus}\quad
\bigl((a_1 b_1) - (a_2 b_2)\bigr)
\bigl((a_2 b_1) + (a_1 b_2)\bigr)
\end{verbatim}
The third kind of situation is a slightly oversize object in running
text, such as $\left\lvert\frac{b'}{d'}\right\rvert$ where the
delimiters produced by \cn{left} and \cn{right} cause too much line
spreading. In that case \ncn{bigl} and \ncn{bigr}\index{big@\cn{big},
\cn{Big}, \cn{bigg}, \dots\ delimiters} can be used to produce
delimiters that are slightly larger than the base size but still able to
fit within the normal line spacing:
$\bigl\lvert\frac{b'}{d'}\bigr\rvert$.

In ordinary \latex/ \ncn{big}, \ncn{bigg}, \ncn{Big}, and \ncn{Bigg}
delimiters aren't scaled properly over the full range of \latex/ font
sizes.  With the \nipkg{amsmath} package they are.%
\index{delimiters!fixed size|)}

\subsection{Vertical bar notations}

The \nipkg{amsmath} package provides commands \cn{lvert}, \cn{rvert},
\cn{lVert}, \cn{rVert} (compare \cn{langle}, \cn{rangle}) to address the
%problem of overloading for the vert bar character \qc{\|}. This
problem of overloading for the vert bar character \qcvert. This
character is currently used in \latex/ documents to represent a wide
variety of mathematical objects: the `divides' relation in a
number-theory expression like $p\vert q$, or the absolute-value
operation $\lvert z\rvert$, or the `such that' condition in set
notation, or the `evaluated at' notation $f_\zeta(t)\bigr\rvert_{t=0}$.
The multiplicity of uses in itself is not so bad; what is bad, however,
is that fact that not all of the uses take the same typographical
treatment, and that the complex discriminatory powers of a knowledgeable
reader cannot be replicated in computer processing of mathematical
documents. It is recommended therefore that there should be a one-to-one
correspondence in any given document between the vert bar character
%\qc{\|} and a selected mathematical notation, and similarly for the
\qcvert\ and a selected mathematical notation, and similarly for the
%double-bar command \ncn{\|}\index{"|@\verb"*+"\"|+}. This immediately
double-bar command \cnvert. This immediately
rules out the use of \qc{|}
%and \ncn{\|}\index{"|@\verb"*+"\"|+} for delimiters, because left and right
and \cnvert\ for delimiters, because left and right
delimiters are distinct usages that do not relate in the same way to
adjacent symbols; recommended practice is therefore to define suitable
commands in the document preamble for any paired-delimiter use of vert
bar symbols:
\begin{verbatim}
\providecommand{\abs}[1]{\lvert#1\rvert}
\providecommand{\norm}[1]{\lVert#1\rVert}
\end{verbatim}
whereupon the document would contain |\abs{z}| to produce $\lvert
z\rvert$ and |\norm{v}| to produce $\lVert v\rVert$.
The \pkg{mathtools} provides the command \cn{DeclarePairedDelimiter}
for defining |\abs|-like macros with scaling delimiters.
\index{delimiters|)}

%%%%%%%%%%%%%%%%%%%%%%%%%%%%%%%%%%%%%%%%%%%%%%%%%%%%%%%%%%%%%%%%%%%%%%%%

\chapter{Operator names}

\section{Defining new operator names}\label{s:opname}

Math functions\index{operator names}\relax \index{function
names|see{operator names}} such as $\log$, $\sin$, and $\lim$ are
traditionally typeset in roman type to make them visually more distinct
from one-letter math variables, which are set in math italic. The more
common ones have predefined names, \cn{log}, \cn{sin}, \cn{lim}, and so
forth, but new ones come up all the time in mathematical papers, so the
\nipkg{amsmath} package provides a general mechanism for defining new
`operator names'. To define a math function \ncn{xxx} to work like
\cn{sin}, you write
\begin{verbatim}
\DeclareMathOperator{\xxx}{xxx}
\end{verbatim}
whereupon ensuing uses of \ncn{xxx} will produce {\upshape xxx} in the
proper font and automatically add proper spacing\index{horizontal
spacing!around operator names} on either side when necessary, so that you
get $A\xxx B$ instead of $A\mathrm{xxx}B$. In the second argument of
\cn{DeclareMathOperator} (the name text), a pseudo-text mode prevails:
the hyphen character \qc{\-} will print as a text hyphen rather than a
minus sign and an asterisk \qc{\*} will print as a raised text asterisk
instead of a centered math star. (Compare
\textit{a}-\textit{b}*\textit{c} and $a-b*c$.) But otherwise the name
text is printed in math mode, so that you can use, e.g., subscripts and
superscripts there.

%%%%xxx%%%\enlargethispage{1\baselineskip}
If the new operator should have subscripts and superscripts placed in
`limits' position above and below as with $\lim$, $\sup$, or $\max$, use
the \qc{\*} form of the \cn{DeclareMathOperator} command:
\begin{verbatim}
\DeclareMathOperator*{\Lim}{Lim}
\end{verbatim}
See also the discussion of subscript placement in
Section~\ref{subplace}.

\goodbreak

The following operator names are predefined:

\vspace{-1\baselineskip}
\begin{ctab}{@{}rlrlrlrl}
\cn{arccos}& $\arccos$ &\cn{deg}& $\deg$ &      \cn{lg}& $\lg$ &        \cn{projlim}& $\projlim$\\
\cn{arcsin}& $\arcsin$ &\cn{det}& $\det$ &      \cn{lim}& $\lim$ &      \cn{sec}& $\sec$\\
\cn{arctan}& $\arctan$ &\cn{dim}& $\dim$ &      \cn{liminf}& $\liminf$ &\cn{sin}& $\sin$\\
\cn{arg}& $\arg$ &      \cn{exp}& $\exp$ &      \cn{limsup}& $\limsup$ &\cn{sinh}& $\sinh$\\
\cn{cos}& $\cos$ &      \cn{gcd}& $\gcd$ &      \cn{ln}& $\ln$ &        \cn{sup}& $\sup$\\
\cn{cosh}& $\cosh$ &    \cn{hom}& $\hom$ &      \cn{log}& $\log$ &      \cn{tan}& $\tan$\\
\cn{cot}& $\cot$ &      \cn{inf}& $\inf$ &      \cn{max}& $\max$ &      \cn{tanh}& $\tanh$\\
\cn{coth}& $\coth$ &    \cn{injlim}& $\injlim$ &\cn{min}& $\min$\\
\cn{csc}& $\csc$ &      \cn{ker}& $\ker$ &      \cn{Pr}& $\Pr$
\end{ctab}
\par\nobreak
\vspace{-1.2\baselineskip}
\begin{ctab}{rlrl}
\cn{varinjlim}&  $\displaystyle\varinjlim$&
\cn{varliminf}&  $\displaystyle\varliminf$\\
\cn{varprojlim}& $\displaystyle\varprojlim$&
\cn{varlimsup}&  $\displaystyle\varlimsup$
\end{ctab}

There is also a command \cn{operatorname} such that using
\begin{verbatim}
\operatorname{abc}
\end{verbatim}
in a math formula is equivalent to a use of \ncn{abc} defined by
\cn{DeclareMathOperator}. This may be occasionally useful for
constructing more complex notation or other purposes. (Use the variant
\cn{operatorname*} to get limits.)

%% ------------------------------------------------------------------ %%

\section{\cn{mod} and its relatives}

Commands \cn{mod}, \cn{bmod}, \cn{pmod}, \cn{pod} are provided to deal
with the special spacing conventions of \qq{mod} notation. \cn{bmod} and
\cn{pmod} are available in \latex/, but with the \nipkg{amsmath} package
the spacing of \cn{pmod} will adjust to a smaller value if it's used in
a non-display-mode formula. \cn{mod} and \cn{pod} are variants of
\cn{pmod} preferred by some authors; \cn{mod} omits the parentheses,
whereas \cn{pod} omits the \qq{mod} and retains the parentheses.
\begin{equation}
\gcd(n,m\bmod n);\quad x\equiv y\pmod b;
\quad x\equiv y\mod c;\quad x\equiv y\pod d
\end{equation}
\begin{verbatim}
\gcd(n,m\bmod n);\quad x\equiv y\pmod b;
\quad x\equiv y\mod c;\quad x\equiv y\pod d
\end{verbatim}

%%%%%%%%%%%%%%%%%%%%%%%%%%%%%%%%%%%%%%%%%%%%%%%%%%%%%%%%%%%%%%%%%%%%%%%%
\enlargethispage{1\baselineskip}

\chapter{The \cn{text} command}\label{text}

The main use of the command \cn{text} is for words or
phrases\index{text fragments inside math} in a display. It is very
similar to the \latex/ command \cn{mbox} in its effects, but has a
couple of advantages. If you want a word or phrase of text in a
subscript, you can type |..._{\textrm{word or phrase}}|, which is slightly
easier than the \cn{mbox} equivalent: |..._{\mbox{\rmfamily\scriptsize| |word|
|or| |phrase}}|. Note that the standard \cn{textrm} command will use the
\nipkg{amsmath} \cn{text} definition, but ensure the \verb+\rmfamily+
font is used.
\begin{equation}
f_{[x_{i-1},x_i]} \text{ is monotonic,}
\quad i = 1,\dots,c+1
\end{equation}
\begin{verbatim}
f_{[x_{i-1},x_i]} \text{ is monotonic,}
\quad i = 1,\dots,c+1
\end{verbatim}

The font used for \cn{text} is the same as that of the surrounding
environment; i.e., within a theorem, the contents of \cn{text} will
be set in italic.

If a math expression is included in a \cn{text} string, it must be
explicitly marked as math (|$...$|).
\[
\partial_s f(x) = \frac{\partial}{\partial x_0} f(x)\quad
  \text{for $x= x_0 + I x_1$.}
\]
\begin{verbatim}
\partial_s f(x) = \frac{\partial}{\partial x_0} f(x)\quad
  \text{for $x= x_0 + I x_1$.}
\end{verbatim}

Function names should \emph{not} be entered as \cn{text}.  Instead,
use \cn{mathrm} or \cn{DeclareMathOperator} as appropriate.  These are
fixed entities that should not change depending on outside content
(such as appearing within a theorem that is set in italic), and in the
case of declared operators, proper spacing\index{horizontal
spacing!around operator names} is applied automatically.

%%%%%%%%%%%%%%%%%%%%%%%%%%%%%%%%%%%%%%%%%%%%%%%%%%%%%%%%%%%%%%%%%%%%%%%%

\end{document}
